\chapter{마지막 산란과 바리오제네시스}

\section{멀리 내다본다는 것은 과거를 본다는 뜻이다}

거리가 다르면 우리는 우주를 서로 다른 시점의 모습으로 본다. 가까운 곳에서는 과거를 되돌아보는 시간이 그다지 중요하지 않다. 일 년 전이나 백 년 전의 우주는 지금의 우주와 크게 다르지 않았다. 그러나 충분히 먼 곳에서는 더 젊은 우주를 본다. 젊은 우주는 밀도가 높고, 밀도가 높은 우주는 더 뜨겁다.

충분히 먼 거리에서는 우주의 온도가 태양 표면의 온도와 비슷했던 시기가 보인다. 광자들이 이곳으로 오는 동안 적색편이되었기 때문에 우리에게는 그렇게 보이지 않는다. 대략 1000의 적색편이 인자는 관측 온도가 같은 비율만큼 낮아졌다는 뜻이다:
\[
T_{\mathrm{obs}} \sim \frac{T_{\mathrm{emit}}}{1+z},
\qquad
1+z \sim 10^3.
\]\footnote{이 식은 00:02:20--00:02:39에서 말로 설명한 적색편이 논증을 간결하게 편집해 옮긴 것이다.}

그 지점에서 우주는 불투명해졌다. 우주는 주로 수소로 이루어져 있었고, 헬륨도 조금 있었다. 온도가 수천 도보다 높으면 수소는 이온화된다. 따라서 그 표면 너머에는 중성 원자가 아니라 이온, 주로 자유롭게 움직이는 양성자와 전자가 있다. 그곳의 광자는 에너지도 더 높아 기본적으로 X선에 해당한다. 원자가 우연히 생기더라도 충돌이나 광자에 의해 곧 다시 이온화된다.

이온화된 기체는 전기 전도체이며, 전도체는 복사가 아무런 방해 없이 통과하게 두지 않는다. 전하들은 전기장에 반응해 앞뒤로 흔들리고 복사를 다시 방출한다. 그래서 그 거리 너머의 우주는 우리에게 불투명하다. 마지막 산란면은 태양의 가시 표면과 비슷하다. 물질로 된 벽이 아니라, 매질이 불투명한 상태에서 투명한 상태로 꽤 갑작스럽게 바뀌는 곳이다.

그곳의 입자 대부분은 광자이며, 그 수는 \(10^8\)배 정도 많다. 오늘날에도 마찬가지지만, 현재 광자의 에너지는 너무 낮아서 현재 에너지 밀도에서 중요한 부분을 차지하지 않는다. 온도가 수소의 이온화 온도 아래로 떨어지면 우주는 투명해져 우리는 그 너머를 볼 수 있다.

직접 볼 수 있는 끝은 그곳이다. 다음 질문은 알려진 원자·핵·입자 물리학으로 과거를 얼마나 더 거슬러 올라갈 수 있느냐는 것이다.

\section{더 먼 과거로}

마지막 산란면 바로 안쪽에는 광자 외에 주로 양성자와 전자가 있고, 헬륨 원자핵과 중양자는 조금 있다. 반양성자와 양전자는 사실상 없다. 하지만 더 거슬러 올라가면 온도가 결국 두 광자가 충돌해 전자와 양전자를 만들 수 있을 만큼 높아진다:
\[
\gamma + \gamma \leftrightarrow e^- + e^+.
\]

기본적인 문턱 조건은 광자들이 쌍의 정지 에너지를 적어도 공급해야 한다는 것이다:
\[
E \ge 2m_e c^2.
\]
그 문턱 아래에서는 보통의 빛으로는 그런 반응이 일어나지 않는다. 그 위에서는 전자--양전자 쌍이 생성될 수 있고, 역반응도 일어날 수 있다. 그 온도의 뜨거운 물질을 담은 솥이 있다면, 그 안에서는 보통 물질의 화학 반응과 같은 역할을 하는 입자 반응들이 일어날 것이다.

문턱을 조금 넘은 정도라면 전자와 양전자보다 광자가 여전히 더 많다. 전자의 정지 질량을 만들려면 큰 에너지가 든다. 하지만 온도가 문턱보다 훨씬 높아지면 전자의 정지 에너지보다 운동 에너지가 더 중요해지고, 밀도도 서로 비슷해진다. 그러면 전자, 양전자, 광자의 밀도가 모두 대략 같아지고, 이들이 사방에서 정신없이 충돌한다.

\begin{figure}[ht]
\centering
\includegraphics[width=0.62\textwidth]{lecture_06_figure_02.png}
\caption{전자와 양성자의 질량 척도를 비교한 쌍생성 문턱 개략도.}
\lectureframe{00:18:47}
\end{figure}

다음 질문은 양성자에게 무슨 일이 일어나느냐는 것이다. 양성자는 이미 존재하고 그때쯤에는 밀도도 매우 높지만, 양성자--반양성자 쌍을 풍부하게 만들 만큼의 열에너지는 아직 없다. 이유는 간단하다:
\[
m_p \approx 2000\,m_e,
\qquad
E \ge 2m_p c^2
\]
이 문턱이 훨씬 더 높기 때문이다.\editorialnote{강의 영상에서는 아래쪽 질량 표기가 위쪽 \(2m_e c^2\)보다 선명하지 않지만, 말로 설명한 비교는 전자와 양성자의 생성 문턱 사이의 비교임이 분명하다.}

이는 반입자가 상당한 양으로 존재하지 않는 항성 내부보다도 훨씬 뜨겁다.

하지만 이것조차 잠정적인 설명일 뿐이다. 양성자와 반양성자가 자유롭게 만들어질 만큼 온도가 높을 때면 양성자를 쿼크로 부술 만큼의 온도이기도 하다. 그러므로 정말 뜨거운 수프는 전자, 양전자, 쿼크, 반쿼크, 광자, 중성미자, 글루온과 다른 종류의 입자들로 이루어진 수프로 설명하는 편이 낫다. 칠판에 적힌 간단한 개략 반응은 다음과 같다:
\[
e^- + e^+ \to \gamma \to q + \bar q.
\]

\begin{classroomqa}
\lecturetimestamp{00:20:05}
\audiencequestion{중간자인가요, 아니면 쿼크 이전 단계의 입자인가요?}
\lecturerresponse{중간자는 쿼크로 이루어져 있습니다. 양성자를 부술 만큼 온도가 높으면 중간자도 부서집니다.}
\end{classroomqa}

지금까지의 추론은 시간을 거슬러 올라가는 방향이었다. 이제 방향을 바꾸어 뜨겁고 대칭적인 수프가 식을 때 무슨 일이 일어나는지 물어보자.

\section{전기적 중성과 대칭적인 뜨거운 우주}

온도가 충분히 높으면 우주에는 입자와 반입자가 매우 높은 밀도로 존재한다. 전자와 양전자는 쌍으로, 쿼크와 반쿼크도 쌍으로 나타나므로, 각 입자와 그 반입자가 거의 같은 양만큼 존재한다고 보는 것이 자연스럽다.

전기적 중성에 관해서는 조금 신중할 필요가 있다. 오늘날 우주는 전기적으로 중성이다. 기억을 돕기 위해 실제 우주가 아니라 삼차원 구면 같은 닫히고 유계인 우주를 상상해 보자. 먼저 2차원 구면을 모형으로 생각하자. 전기력선은 전하가 있는 곳을 제외하면 시작하거나 끝날 수 없다. 양전하는 전기력선을 내보내고, 음전하는 전기력선을 흡수한다. 닫힌 구면에서는 전기력선이 무한대로 빠져나갈 수 없다. 대신 반대편에서 다시 한곳으로 모인다. 따라서 기억을 돕는 비유로는 닫히고 유계인 우주가 순전하를 가질 수 없다고 볼 수 있다. 실제 우주가 문자 그대로 닫혀 있는지가 요점은 아니다. 요점은 큰 규모에서 양전하와 음전하가 놀라울 만큼 정확하게 상쇄되는 것처럼 보인다는 데 있다.

전기적 중성은 양성자의 수와 반양성자의 수가 같다는 \emph{뜻이 아니다}. 전체 양전하가 전체 음전하와 같다는 뜻이다. 양성자와 양전자, 그리고 양전하를 띤 다른 모든 입자를 합한 전하가 전자와 음전하를 띤 다른 모든 입자를 합한 전하와 균형을 이룬다.

알려진 물리학을 신뢰할 수 있는 한계까지, 어쩌면 그보다 조금 더 거슬러 올라가면, 우주는 입자와 반입자가 풍부하게 존재할 만큼 뜨거웠다. 이름을 댈 수 있는 어떤 가속기 실험이든 초기 우주의 어느 시점에는 그에 대응하는 현상이 있었다. 거의 대칭적인 시작을 가정하는 것이 자연스럽다. 전자와 양전자의 수가 같고, 쿼크와 반쿼크의 수도 같으며, 양성자와 반양성자를 중입자와 반중입자의 대리물로 삼으면,
\[
N_p = N_{\bar p}.
\]

그 추측이 매력적인 이유는 물리 법칙이 입자와 반입자 사이에서 거의 대칭이기 때문이다. 양전자는 전자와 매우 비슷하게 행동한다. 완전히 같지는 않지만 거의 그렇다. 바로 그 근사적인 대칭성이 문제다.

\section{우주 역사를 앞으로 돌리면}

우주가 입자와 반입자가 같거나 거의 같은 양으로 존재하는 뜨거운 상태에서 시작해 식어 간다고 가정해 보자.

온도가 충분히 높을 동안에는 에너지 높은 충돌로 양성자와 반양성자가 만들어질 수 있고, 이들은 소멸해 광자가 될 수도 있다. 하지만 결국 온도가 양성자 생성 문턱 아래로 떨어진다. 그러면 양성자--반양성자 쌍은 여전히 소멸해 광자가 될 수 있지만, 우주가 조금 더 팽창하고 나면 그 광자는 에너지를 충분히 잃어 다시 양성자--반양성자 쌍으로 바뀔 수 없다. 이것이 일방통행이다. 같은 일이 나중에 전자와 양전자에 대해서도 일어나는데, 그때의 온도는 양성자 생성 문턱보다 훨씬 낮지만 항성의 온도보다는 여전히 훨씬 높다.

팽창은 미시적인 반응 속도에 비해 느리므로, 이 과정이 진행되는 동안 우주는 대체로 열평형에 가깝다. 여기서 단열 팽창이란 이런 뜻이다. 더 정확히 말하면 엔트로피가 거의 변하지 않는다는 뜻이다. 뜨거운 입자 수프에서 엔트로피는 대략 1 정도인 인자들을 제외하면 본질적으로 입자 수에 해당한다. 따라서 소멸하는 양성자, 반양성자, 전자, 양전자는 대체로 광자로 바뀐다.

양성자와 반양성자의 수가 정확히 같았다면 둘은 모두 소멸했을 것이다. 전자와 양전자의 경우도 마찬가지였을 것이다. 그러면 광자만 남았을 것이다. 하지만 양성자가 반양성자보다 아주 조금이라도 더 많았다면, 짝을 이룬 입자들은 소멸하고 짝을 이루지 못한 초과분은 남았을 것이다. 구체적인 예를 들면, 양성자가 100개이고 반양성자가 50개라면 50개의 반양성자가 양성자 50개를 찾아 소멸하고 양성자 50개가 남는다. 전기적 중성을 유지하려면 전자도 양전자보다 그와 같은 방식으로 더 많이 남아야 한다.

이로써 현재의 광자 수에서 초기의 비대칭을 추론할 수 있다. 오늘날 광자의 수는 양성자 또는 전자의 수보다 약 \(10^9\)에서 \(10^{10}\)배 많다:
\[
\frac{N_\gamma}{N_p} \sim 10^9\text{--}10^{10}.
\]
따라서 초기의 뜨거운 플라스마에는 양성자--반양성자 쌍 십억 개당 양성자가 하나 더 있었을 뿐이다. 그 미세한 불균형이 원자가 지금까지 존재하는 유일한 이유다.

\begin{classroomqa}
\lecturetimestamp{00:42:25}
\audiencequestion{냉각 과정에서 아주 작은 확률로 생기는 초과분이라고 통계적으로 생각하면 안 되나요?}
\lecturerresponse{처음에 대칭적인 상태에서 양성자를 하나 더 얻으려면 반양성자는 만들지 않고 양성자만 만들어야 한다. 이와 함께 음전하도 만들어야 하는데, 예를 들면 전자가 그렇다. 알려진 표준 모형에는 그런 저에너지 과정이 없다. 실험으로 검증된 그 에너지 범위에서는 중입자수가 보존되므로, 보통의 냉각으로는 순초과분이 생기지 않는다.}
\end{classroomqa}

전하가 장애물은 아니다. 도식적인 예는 다음과 같다.
\[
\gamma + \gamma \to p + e^-.
\]\footnote{이 표기는 00:44:16--00:46:35에서 말로 제시된 예를 편집상 옮겨 적은 것이다. 그 목적은 전하가 보존되면서 바리온 수는 위반될 수 있다는 점을 분명히 드러내는 데 있다.}
이 반응은 전체적으로 전기적 중성을 만족한다. 그러나 이는 또 다른 경험적 규칙을 위반한다. 검증된 표준 모형에서는 반양성자를 함께 만들지 않고 양성자를 생성할 수 없다. 여기서 관련된 보존량은 바리온 수이다.
\[
B=(N_p+N_n)-(N_{\bar p}+N_{\bar n}).
\]\footnote{이 간결한 공식은 00:45:18--00:45:40에서 제시된 말로 된 정의를 수식으로 나타낸다.}
알려진 에너지 범위에서는 이 양이 보존되므로, 냉각 역사의 후반부인 충분히 검증된 구간에서는 그 비대칭을 만들어 낼 수 없다.

\section{바리오제네시스}

이것이 바리오제네시스라고 불리는 주제다. 우주가 실제로 중입자와 반중입자의 수가 같은 상태에서 시작했다고 가정해 보자. 그 대칭 상태를 관측된 작은 초과분 쪽으로 기울이려면 물리 법칙이 무엇을 허용해야 할까?

칠판에서는 이를 두 개의 기호식으로 간추렸고,
\[
N_p=N_{\bar p},
\qquad
N_p-N_{\bar p},
\]
그 다음에는 첫 번째 식에서 두 번째 식으로 이어지는 조건들을 나열했다.

\begin{figure}[ht]
\centering
\includegraphics[width=0.72\textwidth]{lecture_06_figure_04.png}
\caption{처음에 대칭적인 상태에서 중입자 비대칭을 만들어 내기 위한 조건.}
\lectureframe{01:22:26}
\end{figure}

\editorialnote{영상 속 번호 매긴 목록은 일부만 판독할 수 있다. 아래의 세 항목은 함께 제시된 설명을 따른 것이다.}

첫 번째 조건은 중입자수 위반이다. 중입자수가 정확한 보존 법칙이라면 순중입자수가 0인 상태에서 시작한 우주는 계속 그 상태로 남는다.

두 번째 조건은 입자와 반입자 사이의 편향이다. 중입자수 위반만으로는 충분하지 않다. 자연이 양성자를 만드는 과정과 그에 대응하는 반양성자 생성 과정을 같은 확률로 허용한다면 평균적인 불균형은 여전히 0이다. 두 과정의 반응률이 달라야 한다.

광자를 포함한 전체 입자 수가 \(10^{90}\) 정도라면 통계적 요동의 크기는 \(\sqrt{N}\) 정도이고, 그 결과 생기는 초과분은 겨우 \(10^{45}\) 정도다. 우주에 있는 약 \(10^{80}\)개의 양성자를 설명하기에는 턱없이 부족하다.

세 번째 조건은 우주가 어느 시점에 충분히 비평형 상태에 있었어야 한다는 것이다. 여기서는 CPT 정리를 이용한다. 뜨거운 기체가 영원히 평형 상태에 있었다면, 그 역사를 거꾸로 돌린 영상은 앞으로 돌린 영상과 같아 보인다. 그런 상황에서는 전하 켤레 변환, 공간 반전, 시간 반전을 결합한 정확한 대칭성으로 인해 입자와 반입자 사이의 대칭성이 사실상 회복된다. 실제 편향이 남도록 하려면 미시적인 시간 척도에서 실제 역사가 앞으로 진행될 때와 뒤로 진행될 때 다르게 보여야 한다. 우주의 충분히 초기에는 팽창이 그럴 만큼 빨랐다.

따라서 세 가지 조건은 다음과 같다:
\begin{enumerate}
\item 중입자수 위반;
\item 입자와 반입자 사이 대칭성의 위반, 더 정확히는 CP 위반;
\item 열평형에서 벗어난 단계.
\end{enumerate}

\begin{classroomqa}
\lecturetimestamp{01:09:08}
\audiencequestion{이 조건들을 CPT와 대응시키면, 조건 2가 실제로 CP이고 조건 3가 실제로 시간 방향에 관한 문제인가요?}
\lecturerresponse{그렇습니다. 조건 2는 정확히 CP 위반이라고 해야 합니다. 조건 3는 실제 역사가 비평형 상태였다는 사실입니다. 이는 별도의 보존 법칙도 아니고 중입자수 위반과 같은 것도 아닙니다.}
\end{classroomqa}

이 세 조건이 갖춰지면 양성자와 반양성자의 수가 같은 상태에서 시작한 우주라도 어느 한쪽으로 기울어질 수 있다. 아직 확실하지 않은 점은 자연에서 실제로 중입자수가 위반되는지, 그렇다면 그 방식이 무엇인지다.

\section{알려진 것과 알려지지 않은 것}

이 이야기에서 아직 밝혀지지 않은 부분은 중입자수 위반이다. 자연에서 다음을 변화시키는 반응을 직접 관측한 적은 없다:
\[
(N_p+N_n)-(N_{\bar p}+N_{\bar n}).
\]
관련 반응이 실험실에서 직접 검증된 어떤 온도보다 훨씬 높은 온도를 필요로 하거나, 낮은 에너지에서는 발생 확률이 엄청나게 작다는 뜻일 수도 있다. 하지만 현재로서는 여전히 물음표로 남아 있다.

그래서 양성자 붕괴가 이 논의에서 매우 중요하다. 중입자수가 정확히 보존되지 않는다면 양성자는 원칙적으로 불안정해야 한다. 하지만 양성자는 놀라울 만큼 안정하다. 여기서 제시하는 수명 하한은 대략 다음과 같다:
\[
\tau_p \gtrsim 10^{34}\ \text{years}.
\]
이는 빛 검출기를 설치한 매우 큰 부피의 물을 이용해 지하 깊은 곳에서 수행한 대규모 실험에서 얻은 값이다. 이 결과가 중입자수 위반을 반증하는 것은 아니다. 그런 위반이 있더라도 낮은 에너지에서는 극도로 억제되어야 한다는 뜻일 뿐이다.

알려진 물리학에 관한 별도의 언급도 논의에 등장한다. \(100\)에서 \(250\,\mathrm{GeV}\) 사이에서는 약한 상호작용이 전자기 상호작용과 훨씬 더 비슷해진다. 이는 구체적인 모형에서는 중요할 수 있지만, 현재 논의에서 극적인 전환점은 아니다.

암흑 물질도 기본적인 논리를 크게 바꾸지는 않는다. 대부분의 이론에서 암흑 물질은 보통 물질과 매우 약하게 상호작용하므로, 이 질문에서는 사실상 방관자다.

\begin{classroomqa}
\lecturetimestamp{01:24:49}
\audiencequestion{중력은 이 모든 과정에 어떤 영향을 주나요?}
\lecturerresponse{주로 두 가지 방식으로 영향을 줍니다. 첫째, 우주 팽창이 없었다면 비평형 단계도 없었을 것입니다. 둘째, 블랙홀 증발은 중입자수가 영원히 정확한 보존량일 수 없다고 의심할 강력한 이유를 제공합니다. 그 밖에는 입자 사이의 일반적인 중력 상호작용이 여기서 중요한 영향을 미치기에는 너무 약합니다.}
\end{classroomqa}

그러면 전하에 관한 후속 질문이 생긴다. 같은 블랙홀 논증이 전하 보존도 위협할까? 답은 아니다. 전하를 띤 블랙홀은 바깥으로 뻗는 전기력선을 가지며, 그 전기력선은 전하를 기록한다. 그런 블랙홀은 완전히 증발하지 않고 극한 잔해에 가까워진다. 따라서 블랙홀 논증이 위협하는 것은 중입자수이지 전하가 아니다.

또 다른 질문은 초기의 밀도 높은 쿼크 수프를 그저 블랙홀로 보아야 하느냐는 것이다. 답은 아니다. 요점은 정지해 있는 고밀도 덩어리가 아니다. 우주는 팽창하고 있었다. 블랙홀 이야기는 중입자수의 운명을 살펴보는 사고실험이지, 초기 우주가 실제로 갇힌 영역이었다는 문자 그대로의 설명이 아니다.

\begin{classroomqa}
\lecturetimestamp{01:31:14}
\audiencequestion{그런 반응들을 포함하도록 표준 모형을 어떻게 수정해야 하는지 알 수 있을까요?}
\lecturerresponse{문제는 그 반대입니다. 중입자수 위반을 지나치게 많이 추가해 양성자가 너무 빨리 붕괴하게 만들기는 쉽습니다. 원하는 초기 우주 과정을 허용하면서도 양성자가 지금처럼 안정하도록 유지하는 것이 진짜 문제입니다.}
\end{classroomqa}

\section{엔트로피와 뒤이어 나온 중성자 질문}

이제 다시 수를 따져 보자. 단열 팽창하는 우주에서는 엔트로피가 대략 보존된다. 뜨거운 기본 입자 수프에서 엔트로피는 1 정도인 인자들을 제외하면 본질적으로 입자 수에 해당한다. 그래서 이 논의에서는 후기에 아무것도 없는 곳에서 엄청난 양의 광자가 생겨날 필요가 없다. 소멸 과정은 대부분 수프의 구성 입자들을 복사로 바꾸어 놓는다.

문제의 이전 버전은 상황을 거꾸로 놓고 보았다. 광자에 비해 양성자가 왜 그렇게 적게 살아남았는지를 묻는 대신, 질문을 반대로 제기했다. 양성자와 비교해 광자는 왜 그렇게 많은가? 그러나 더 날카로운 질문은 우주가 왜 그렇게 늦은 시기에 그토록 많은 엔트로피를 만들어 냈는가가 아니다. 애초에 입자와 반입자 사이의 비대칭은 왜 그렇게 미미했는가 하는 것이다. 양성자--반양성자 쌍 십억 개당 양성자가 겨우 하나 더 있었을 뿐이다.

자유 중성자는 붕괴하지만, 그 과정은 중입자수를 위반\emph{하지 않는다}:
\[
n \to p + e^- + \bar\nu.
\]
중성자는 양성자와 같은 중입자수를 가진다.

자유 중성자에게 허용되는 붕괴 에너지는 매우 작다. 중성자의 질량이 최종 상태의 질량보다 아주 조금 클 뿐이기 때문이다. 이것이 기본 입자 물리학의 기준으로 볼 때 붕괴가 매우 느린 이유 중 하나다. 중성자가 원자핵 안에 결합되어 있으면 결합 에너지가 그 작은 가용 에너지를 완전히 없앨 수 있고, 그러면 결합된 중성자는 안정할 수 있다.

\begin{classroomqa}
\lecturetimestamp{01:43:03}
\audiencequestion{중성자는 우주에서 사라지기만 하나요, 아니면 생성되기도 하나요?}
\lecturerresponse{첫 번째 답은 자유 중성자가 붕괴로 사라진다는 것입니다. 하지만 그것이 전부는 아닙니다. 별에서 일어나는 핵반응과 중성자별의 형성을 비롯해 실제 중성자 생성원도 있습니다. 그러므로 전체 중성자 수가 한 방향으로만 줄어드는 것은 아닙니다.}
\end{classroomqa}